\documentclass[14 pt]{extarticle}

	\usepackage[frenchb]{babel}
	\usepackage[utf8]{inputenc}  
	\usepackage[T1]{fontenc}
	\usepackage{amssymb}
	\usepackage[mathscr]{euscript}
	\usepackage{stmaryrd}
	\usepackage{amsmath}
	\usepackage{tikz}
	\usepackage[all,cmtip]{xy}
	\usepackage{amsthm}
	\usepackage{varioref}
	\usepackage{geometry}
	\geometry{a4paper}
	\usepackage{lmodern}
	\usepackage{hyperref}
	\usepackage{array}
	 \usepackage{fancyhdr}
\renewcommand{\theenumi}{\alph{enumi})}
	\pagestyle{fancy}
	\theoremstyle{plain}
	\fancyfoot[C]{} 
	\fancyhead[L]{Contrôle}
	\fancyhead[R]{4 juin 2025}\geometry{
 a4paper,
 total={170mm,257mm},
 left=20mm,
 top=20mm,
 }
	
	
	\title{Contrôle proportionnalité}
	\date{}
	\begin{document}

 
 \subsection*{Calculs (6 points)}
 Effectuer en détaillant vos calculs les opérations suivantes : 
 \begin{enumerate}
 \item $3 - 2 \times 4$
 \item $-1 + 5$
 \item $ - 12 -3$
 \item $14- (3-5)$
 \item $9 - ( - 5 +2)$
 \item $12 - 4 \times (5 - (-1))$
 \end{enumerate}
 \subsection*{Exercice 1 (4 points)}

 Reconnaître parmi les tableaux suivants ceux qui sont des tableaux de proportionnalité. Justifier vos réponses.
 \begin{enumerate}
 \item $\renewcommand{\arraystretch}{1}
\begin{array}{|c|c|c|}
\hline
\phantom{000}21\phantom{000} & \phantom{000}49\phantom{000} & \phantom{000}63\phantom{000}\\
\hline
3 & 7 & 9\\
\hline
\end{array}
\renewcommand{\arraystretch}{1}$

\item $\renewcommand{\arraystretch}{1}
\begin{array}{|c|c|c|}
\hline
\phantom{000}72\phantom{000} & \phantom{000}54\phantom{000} & \phantom{000}18\phantom{000}\\
\hline
60 & 45  & 15\\
\hline
\end{array}
\renewcommand{\arraystretch}{1}$

 \item $\renewcommand{\arraystretch}{1}
\begin{array}{|c|c|c|}
\hline
\phantom{000}7\phantom{000} & \phantom{000}46 \phantom{000} & \phantom{000}6\phantom{000}\\
\hline
4 & 8 & 3\\
\hline
\end{array}
\renewcommand{\arraystretch}{1}$



\end{enumerate}

\subsection*{Exercice 2 (6 points)}

Recopier et remplir les tableaux de proportionnalité suivants. 

\begin{enumerate}
\item $\renewcommand{\arraystretch}{1}
\begin{array}{|c|c|c|c|c|}
\hline
\phantom{000}20 \phantom{000} &\phantom{000}5\phantom{000} &  \phantom{000}35\phantom{000} &   \phantom{000}55\phantom{000} & \phantom{000} \phantom{000} \\
\hline
120 & \phantom{700} & &   & 720 \\
\hline
\end{array}
\renewcommand{\arraystretch}{1}$

\item $\renewcommand{\arraystretch}{1}
\begin{array}{|c|c|c|c|c|}
\hline
\phantom{000}6\phantom{000}   & \phantom{000} \ \phantom{000} & \phantom{000}28\phantom{000}& \phantom{000}\phantom{000}   \\
\hline
45 &  90& &1650\\
\hline
\end{array}
\renewcommand{\arraystretch}{1}$
\end{enumerate}

\subsection*{Exercice 3 (4 points)}

Un train réalise le trajet de Paris à Lyon en $2$ heures. 

a) Sachant que la distance entre Paris et Lyon est de $480$ 
kilomètres, calculer la vitesse moyenne du train sur le trajet, en kilomètres par heure, puis en kilomètres par minute. 

b) La distance entre Paris et Marseille est de $776$ kilomètres. En 
supposant que la vitesse moyenne du train n'a pas changé, calculer la durée d'un trajet Paris-Marseille. (Exprimer le résultat sous la forme \_ heures et \_ minutes.)

 
 \newpage
 
 
 \subsection*{Calculs (6 points)}
 Effectuer en détaillant vos calculs les opérations suivantes : 
 \begin{enumerate}
 \item $3 - 4 \times 3$
 \item $-1 + 8$
 \item $ - 11 -3$
 \item $13- (3-5)$
 \item $9 - ( - 6 +2)$
 \item $12 - 3 \times (5 - (-1))$
 \end{enumerate}
 \subsection*{Exercice 1 (4 points)}

 Reconnaître parmi les tableaux suivants ceux qui sont des tableaux de proportionnalité. Justifier vos réponses.
 \begin{enumerate}
 \item $\renewcommand{\arraystretch}{1}
\begin{array}{|c|c|c|}
\hline
\phantom{000}21\phantom{000} & \phantom{000}56\phantom{000} & \phantom{000}63\phantom{000}\\
\hline
3 & 8 & 9\\
\hline
\end{array}
\renewcommand{\arraystretch}{1}$

\item $\renewcommand{\arraystretch}{1}
\begin{array}{|c|c|c|}
\hline
\phantom{000}72\phantom{000} & \phantom{000}54\phantom{000} & \phantom{000}18\phantom{000}\\
\hline
60 & 45  & 15\\
\hline
\end{array}
\renewcommand{\arraystretch}{1}$

 \item $\renewcommand{\arraystretch}{1}
\begin{array}{|c|c|c|}
\hline
\phantom{000}7\phantom{000} & \phantom{000}46 \phantom{000} & \phantom{000}6\phantom{000}\\
\hline
4 & 8 & 3\\
\hline
\end{array}
\renewcommand{\arraystretch}{1}$



\end{enumerate}

\subsection*{Exercice 2 (6 points)}

Recopier et remplir les tableaux de proportionnalité suivants. 

\begin{enumerate}
\item $\renewcommand{\arraystretch}{1}
\begin{array}{|c|c|c|c|c|}
\hline
\phantom{000}30 \phantom{000} &\phantom{000}5\phantom{000} &  \phantom{000}25\phantom{000} &   \phantom{000}55\phantom{000} & \phantom{000} \phantom{000} \\
\hline
120 & & \phantom{700} &    & 720 \\
\hline
\end{array}
\renewcommand{\arraystretch}{1}$

\item $\renewcommand{\arraystretch}{1}
\begin{array}{|c|c|c|c|}
\hline
\phantom{000}6\phantom{000}  & \phantom{000} \ \phantom{000} & \phantom{000}28\phantom{000}& \phantom{000}\phantom{000}   \\
\hline
45  & 90& &1650\\
\hline
\end{array}
\renewcommand{\arraystretch}{1}$
\end{enumerate}

\subsection*{Exercice 3 (4 points)}

Un train réalise le trajet de Paris à Lyon en $2$ heures. 

a) Sachant que la distance entre Paris et Lyon est de $480$ 
kilomètres, calculer la vitesse moyenne du train sur le trajet, en kilomètres par heure, puis en kilomètres par minute. 

b) La distance entre Paris et Marseille est de $776$ kilomètres. En 
supposant que la vitesse moyenne du train n'a pas changé, calculer la durée d'un trajet Paris-Marseille. (Exprimer le résultat sous la forme \_ heures et \_ minutes.)

 
 


 	\end{document}
